\section{命名样式与缺失值}
\key{\gradbarsstyle{名称}{参数}} 定义可复用的配置；使用 \key{style=名称} 调用。
样式可以引用已定义的样式，后面的参数覆盖前面的设置。
定义和调用均遵循普通 TeX 分组，不会改变其他单元格。

\begin{example}{示例 6：定义并复用样式}
\gradbarsstyle{accuracy}{
  max=100,theme=teal,
  rounded=2pt,precision=1,
  unit={\%}
}
\gradbar[style=accuracy]{92.7}

\gradbar[style=accuracy,
  target=95]{94.2}
\end{example}

样式在当前作用域内可以重新定义。未知名称会报错；循环引用会触发嵌套深度限制。

\begin{example}{示例 7：零值与缺失值}
\gradbar{0}

\gradbar{NA}

\gradbar[missing text={N/A}]{}
\end{example}

空参数或大写 \key{NA} 表示缺失值，忽略首尾空格。
零值显示数值零；缺失值只显示轨道和缺失标签，默认标签为短横线。
宏包不会把其他文字或拼写错误自动当成缺失值。

\note{缺失值不绘制数值填充和误差线，也不附加单位或百分号。
目标线和区间背景仍保留，表示这一列的参考尺度。
\key{text} 可以覆盖缺失标签，\key{label=none} 可以隐藏它。}
